%------------------------------------------------------------------------------%
\section{Série Alternada e Teste de Leibniz}

Nas seções anteriores apresentamos vários testes que verificam se uma
série de termos não negativos é convergente. Agora vamos discutir
séries que possuem termos positivos e negativos. Nesta seção
estudaremos o caso particular onde os sinais se alternam produzindo o que
chamamos de série alternada. Na próxima seção apresentamos os conceitos de
convergência absoluta e condicional que vamos usar para os outros casos.

%------------------------------------------------------------------------------%
\begin{definicao}{Série Alternada}{def:serie-alternada}
  Uma série com a forma
  \[
    \sum_{n=1}^\infty (-1)^{n+1}\,a_n
    = a_1 - a_2 + a_3 - a_4 + \cdots
  \]
  onde $a_n>0$ para todo $n$ é chamada de
\emph{Série Alternada}\index{Série!alternada}.
\end{definicao}
%------------------------------------------------------------------------------%

Para esse caso particular de série temos um teste de convergência bastante
simples, o Teste de Leibniz~\ref{teo:Teste-Leibniz}, e também uma forma de
estimar o erro ao aproximarmos a soma da série
por uma soma parcial~\ref{pro:erro-serie-alternada}.

%------------------------------------------------------------------------------%
\begin{teorema}{Teste de Leibniz}{teo:Teste-Leibniz}
  Seja\index{Teste!Leibniz}
  \[
    \sum_{n=1}^\infty (-1)^{n+1}\,a_n
  \]
  uma Série Alternada, Definição~\ref{def:serie-alternada}, se
  \begin{enumerate}
      \item $a_n>0$,
      \item $a_n\geq a_{n+1}$,
      \item $a_n\to 0$
  \end{enumerate}
  para todo $n \geq N$, então a série converge.
\end{teorema}
%------------------------------------------------------------------------------%
\begin{proof}
Vamos assumir que $N=1$, isso é, as condições do teste valem para todos os
termos da série. Escolhendo $n$ par e escrevendo $n=2m$, temos que
a soma dos $n$ primeiros termos é
\[
  S_n = S_{2m} = a_1 - a_2 + a_3 - a_4 + a_5 - a_6 + \cdots + a_{2m-1} - a_{2m}    \,.
\]
Como essa é uma soma finita podemos agrupar os termos dois a dois
\[
  S_{2m} = (a_1 - a_2) + (a_3 - a_4) + (a_5 - a_6) + \cdots + (a_{2m-1} - a_{2m})    \,.
\]
Cada termo entre parenteses é positivo ou zero, pois $a_k\geq a_{k+1}$,
então podemos concluir que $S_{2m+2}\geq S_{2m}$ para todo $m$, ou seja,
a sequência $S_{2m}$ é crescente.
Vamos agora agrupar os termos de $S_{2m}$ de outra forma
\[
  S_{2m} = a_1 - (a_2 - a_3) - (a_4 - a_5) - \cdots -  (a_{2m-2} - a_{2m-1}) - a_{2m}
\]
dessa igualdade concluímos que
\(
  S_{2m} \leq a_1
\).
Como a sequência $S_{2m}$ é crescente e limitada, pelo Teorema~\ref{teo:sequencia-monotona},
concluirmos que ela é convergente e podemos escrever
\begin{equation}
  \label{eq:test_leibniz_par}
  \lim_{m\to \infty} S_{2m} = L
\end{equation}
para algum $L\in\R$.
Considerando agora o caso onde $n$ é ímpar, $n=2m+1$, temos que a soma dos $n$
primeiros termos é
\[
  S_n = S_{2m+1} = S_{2m} + a_{2m+1}
\]
como $a_n\to 0$ temos que
\begin{equation}
  \label{eq:test_leibniz_impar}
  \lim_{m\to \infty} S_{2m+1}
  = \lim_{m\to \infty} S_{2m} + \lim_{m\to \infty} a_{2m+1}
  = L + 0 = L   \,.
\end{equation}
A equação~\eqref{eq:test_leibniz_par} nos informa que a subsequência dos termos pares
tende para~$L$, enquanto que, a equação~\eqref{eq:test_leibniz_impar} garante que os
termos impares convergem para o mesmo valor. Podemos concluir então que
\[
  \sum_{n=1}^\infty (-1)^{n+1}a_n
  = \lim_{n\to \infty} S_n
  = L    \,.
\]
\end{proof}
%------------------------------------------------------------------------------%

Para aplicarmos o Teste de Leibniz precisamos garantir que $a_n\geq a_{n+1}$,
o que pode ser complicado em alguns casos. Porém, se conhecermos uma função $f$
tal que $f(n)=a_n$, podemos verificar se a função é decrescente calculando sua
derivada. Se $f'(x)\leq 0$ para todo $x\leq N$ podemos concluir que
$a_n\geq a_{n+1}$ para todo $n\geq N$.

Um resultado que podemos extrair do Teste de Leibniz é o critério de convergência
da Série-$p$ com termos alternados.

%------------------------------------------------------------------------------%
\begin{proposicao}{Convergência da Série~$p$ Alternada}{prop:convergencia-serie-p-alternada}
  Observe que $a_n=\sfrac{1}{n^p}$ com $p>0$, satisfaz todas as hipóteses do
  Teorema~\ref{teo:Teste-Leibniz},
  então a \emph{Série~$p$ Alternada}\index{Série!$p$ alternada}
  \[
    \sum_{n=1}^\infty \frac{\,(-1)^{n+1}\,}{n^p}
  \]
  é convergente para todo $p>0$.
\end{proposicao}
%------------------------------------------------------------------------------%

Note para $0<p\leq 1$ a série~$p$ não converge,
entretanto a série~$p$ alternada converge.

Apresentamos a seguir alguns exemplos do uso do Teste de Leibniz.

%------------------------------------------------------------------------------%
\begin{example}
  Use o Teste de Leibniz para mostrar que a série a seguir converge
  \[
    \sum_{n=1}^\infty(-1)^{n}\ln\left(1+\frac{1}{n}\right)     \,.
  \]
  \solution
  Seja
  \[
    a_n=\ln\left(1+\frac{1}{n}\right)    \,,
  \]
  note que $a_n>0$ para todo $n\geq 1$, pois
  \(
    1 + \frac{1}{n} > 1
  \),
  note também que se definirmos
  \[
    f(x) = \ln\left(1+\frac{1}{x}\right)
  \]
  temos
  \[
    f'(x)=\frac{-1}{x^2+x}<0 \qquad \forall x\geq 1
  \]
  o que mostra que $a_{n+1}<a_n$.
  Então $a_n$ é decrescente. Além disso, como a função $\ln(x)$ é contínua em $1$
  segue que
  \[
    \lim a_n = \ln\lim \left(1+\frac{1}{n}\right)
             = \ln 1
             = 0     \,.
  \]
  Assim, podemos aplicar o Teste de Leibniz e concluir que a série alternada converge.
\end{example}
%------------------------------------------------------------------------------%

%------------------------------------------------------------------------------%
\begin{example}
  Use o Teste de Leibniz para mostrar que a série a seguir converge
  \[
    \sum_{n=2}^\infty(-1)^{n+1}\frac{4}{(\ln n)^2}   \,.
  \]
  \solution
  Como  $\ln n$ é crescente,
  \[
    a_n = \frac{4}{(\ln n)^2}
  \]
  é decrescente.
  Além disso $a_n>0$ para todo $n\geq 2$ e $a_n\to 0$, pois
  \[
    \left(\ln n\right)^2\to \infty    \,.
  \]
  Dessa forma, podemos aplicar o teste de Leibniz para concluir que a série alternada converge.
\end{example}
%------------------------------------------------------------------------------%

%------------------------------------------------------------------------------%
\begin{example}
  Use o teste de Leibniz para mostrar que a série a seguir converge
  \[
    \sum_{n=1}^\infty(-1)^{n+1}\frac{10^n}{(n+1)!}        \,.
  \]
  \solution
  Seja
  \(
    a_n=\frac{10^n}{(n+1)!}
  \).
  Veja que se $n\geq 8$, então $n+2\geq 10$ e
  \(
    1 \geq \frac{10}{n+2}
  \).
  Assim,
  \begin{align*}
    a_{n+1} & =\frac{10^{n+1}}{(n+2)!}             \\
            & =\frac{10^n}{(n+1)!}\frac{10}{(n+2)} \\
            & \leq \frac{10^n}{(n+1)!}             \\
            & =a_n
  \end{align*}
  ou seja, $a_n$ é decrescente a partir de $n=8$.

  Para verificar que $a_n\to 0$, note que se $n\geq 10$,
  \begin{align*}
    (n+1)! & = 1\cdot 2\cdot 3\cdots n\cdot (n+1) \\
           & \geq 10! \; 10^{n-10}\,(n+1)         \,.
  \end{align*}
  Dessa forma
  \begin{align*}
    0 & <\frac{10^n}{(n+1)!}                    \\
      & \leq \frac{10^n}{10!\;10^{n-10}\,(n+1)} \\
      & =\frac{10^{10}}{10!\,(n+1)}             \,.
  \end{align*}
  Então, pelo Teorema do Confronto~\ref{teo:confronto} temos que $a_n\to 0$.

  Desse modo, todas as hipóteses do teorema são satisfeitas,
  de forma que podemos concluir que essa série alternada converge.
\end{example}
%------------------------------------------------------------------------------%

%------------------------------------------------------------------------------%
%------------------------------------------------------------------------------%
\subsection*{Estimativa de Erro}
%------------------------------------------------------------------------------%
%------------------------------------------------------------------------------%

Assim como no Teste da Integral o Teste de Leibniz também nos fornece uma
estimativa para o erro $R_n$ de uma série que atende às condições do teste.

%------------------------------------------------------------------------------%
\begin{proposicao}{Erro de uma Série Alternada}{pro:erro-serie-alternada}
  Seja%
  \index{Série!estimativa de erro!teste Leibniz}%
  \index{Estimativa de erro!teste Leibniz}
  $a_k>0$ tal que $a_{k+1}\leq a_k$ e que a série alternada converge para $S$
  \[
    \sum_{k=1}^\infty (-1)^{k+1}\,a_k     \,.
  \]
  Nesse caso, o erro cometido ao aproximarmos $S$ por uma soma parcial
  \[
    S_n = \sum_{k=1}^n (-1)^{k+1}\,a_k
  \]
  é o resto da série (Definição~\ref{def:resto-serie}) e está limitado pelo
  primeiro termo não incluído na soma parcial
  \[
    \abs{R_n} = \abs{S - S_n} \leq \abs{a_{n+1} }    \,.
  \]
\end{proposicao}
%------------------------------------------------------------------------------%
\begin{proof}
  Note que, o resto da série
  \[
    R_n = \sum_{k=n+1}^\infty(-1)^{k+1}\,a_k
  \]
  pode ser expandido em
  \[
    R_n =
    \begin{cases}
      \hphantom{-} a_{n+1}-a_{n+2}+a_{n+3}-\cdots, & \;\text{ se } n \text{ é par} \\
      -            a_{n+1}+a_{n+2}-a_{n+3}+\cdots, & \;\text{ se } n \text{ é ímpar}
    \end{cases}
  \]
  Podemos escrever então que
  \[
    \abs{R_n} = \abs{\,a_{n+1}-(a_{n+2}-a_{n+3})-(a_{n+4}-a_{n+5})-\cdots\,}     \,.
  \]
  Agora, usando que todos os termos entre parênteses são positivos,
  pois $a_{n}>a_{n+1}$, temos que
  \[
    \abs{R_n} < \abs{a_{n+1}}
  \]
  o que demonstra o teorema.
\end{proof}
%------------------------------------------------------------------------------%

%------------------------------------------------------------------------------%
